\chapter{Conclusion}
\label{ch:conclusion}

The thesis asked whether choosing demonstrations can help a relatively small model, at 7B, follow them over its own priors under controlled distribution shift. To answer it, we built a synthetic tabular testbed in which covariate shift, spurious reversal and semantic-prior conflict can each be switched on alone; semantic-prior conflict keeps the rule fixed and sets the model's semantic prior against it by naming the features. Every confirmatory claim was then tested against a contrast family frozen before its data was seen.

The main takeaway is that, in our setting, the prior carried by the feature names influenced the model's predictions more than the choice of demonstrations did. For Qwen2.5-7B-Instruct with 8 numeric demonstrations, flipping three feature names cost 0.20 \AUROC, more than any shift of the data, while the hand-designed strategies typically differed from balanced random demonstrations by only 0.025, and none of the selection methods we tried overcame the names. The one clear effect of selection is that neutralising the shortcut fares better under its reversal than under covariate shift. A learned selector, the only method with no weak environment on a stand-in learner, evened out the environments on the instruction-tuned model at the cost of lower performance in distribution and under covariate shift, and was not significantly better than balanced random demonstrations; in an exploratory study it did help the base model, whose priors pull less, without that cost. Inside the model, the logit lens and probes suggest that the answers form in the last few layers, where the names' effect appears; the prior is readable from the start and the demonstrated rule hardly at all, and a single steering direction moves the answers without making them more correct. These tools locate the effect but do not fully explain how the model decides.

Selection chooses evidence; in our experiments, it did not make the instruction-tuned model read evidence it largely ignored. Within the scale we could study, this appears to be the binding constraint, although our null results bound how much selection helped here rather than show that it cannot help. The testbed we have constructed here, including the pre-registered families and the diagnosis, gives future work, at a larger scale and with the model's own training signal, a controlled place to find out when it stops being a constraint.
